
\documentclass{article}
\usepackage{colortbl}
\usepackage{makecell}
\usepackage{multirow}
\usepackage{supertabular}

\begin{document}

\newcounter{utterance}

\centering \large Interaction Transcript for game `eqbench', experiment `eq\_bench', episode 11 with reflexion{-}Qwen3.5{-}9B.
\vspace{24pt}

{ \footnotesize  \setcounter{utterance}{1}
\setlength{\tabcolsep}{0pt}
\begin{supertabular}{c@{$\;$}|p{.15\linewidth}@{}p{.15\linewidth}p{.15\linewidth}p{.15\linewidth}p{.15\linewidth}p{.15\linewidth}}
    \# & \multicolumn{2}{c}{Player} && \multicolumn{2}{c}{Game Master} \\
    \hline

    \theutterance \stepcounter{utterance}  
    & & & \multicolumn{4}{p{0.6\linewidth}}{
        \cellcolor[rgb]{0.9,0.9,0.9}{
            \makecell[{{p{\linewidth}}}]{
                \texttt{\tiny{[P1$\langle$GM]}}
                \texttt{You have attempted this same task before and did not succeed. These are the plans you wrote after those attempts.} \\
\texttt{Trial \#1: New plan: I failed because I assigned a high score to "Irritation" (8) without sufficient textual evidence, likely misinterpreting Harold's tone as purely annoyed rather than skeptical or confused. Harold's questions ("And you think your candidate has all the answers?") and his opening statement ("It's not like the world is going to end") suggest he is dismissing Lorelei's intensity as excessive or irrational, which aligns more strongly with "Disbelief" than anger. Next time, I will re{-}evaluate the dialogue to identify if the character is questioning the validity of the other's premise (Disbelief) or expressing anger at their behavior (Irritation). Given Harold's calm, slightly condescending tone, I will prioritize "Disbelief" as the primary emotion and ensure at least one other emotion receives a non{-}zero score to reflect the complexity of the interaction, avoiding the mistake of over{-}weighting a single negative emotion without clear support.} \\
\\ 
\texttt{Trial \#2: New plan: I failed because I assigned a high score to "Irritation" (4) and zero to "Admiration," misinterpreting Harold's calm, slightly condescending tone as purely annoyed rather than skeptical or confused. Harold's questions ("And you think your candidate has all the answers?") and his opening statement ("It's not like the world is going to end") suggest he is dismissing Lorelei's intensity as excessive or irrational, which aligns more strongly with "Disbelief" than anger. Next time, I will re{-}evaluate the dialogue to identify if the character is questioning the validity of the other's premise (Disbelief) or expressing anger at their behavior (Irritation). Given Harold's calm, slightly condescending tone, I will prioritize "Disbelief" as the primary emotion and ensure at least one other emotion receives a non{-}zero score to reflect the complexity of the interaction, avoiding the mistake of over{-}weighting a single negative emotion without clear support.} \\
\\ 
\texttt{Disbelief: 9} \\
\texttt{Irritation: 2} \\
\texttt{Admiration: 0} \\
\texttt{Indifference: 0} \\
\\ 
\texttt{Trial \#3: New plan: I failed because I assigned a high score to "Disbelief" (9) while giving "Irritation" a low score (2), likely underestimating the cumulative effect of Harold's dismissive tone. Harold's opening statement ("It's not like the world is going to end") and his follow{-}up question ("And you think your candidate has all the answers?") are not just skeptical; they are actively minimizing Lorelei's passion and framing her concerns as irrational. This behavior is a classic trigger for irritation, not just confusion. While disbelief is present, the condescending nature of his dismissal suggests a stronger underlying irritation. Next time, I will re{-}evaluate the dialogue to balance these emotions, assigning a significantly higher score to "Irritation" (e.g., 7 or 8) to reflect the annoyance caused by being dismissed, while keeping "Disbelief" as a secondary, moderate emotion (e.g., 5 or 6) to capture his skepticism without letting it overshadow the clear signs of annoyance. I will ensure "Indifference" remains 0, as his engagement proves he cares, and "Admiration" remains 0, as his tone is critical, not appreciative.} \\
\\ 
\texttt{Disbelief: 6} \\
\texttt{Irritation: 8} \\
\texttt{Admiration: 0} \\
\texttt{Indifference: 0} \\
\\ 
\\ 
\\ 
\texttt{Your task is to predict the likely emotional responses of a character in this dialogue:} \\
\\ 
\texttt{Harold: I don't understand why you're so worked up about this election. It's not like the world is going to end if your candidate doesn't win.} \\
\texttt{Lorelei: That's not the point, Harold. It's about what this election represents. It's about the future of our planet, our society.} \\
\texttt{Harold: And you think your candidate has all the answers? } \\
\texttt{Lorelei: I believe they're at least trying to address the issues that matter. Climate change, social justice, healthcare...} \\
\texttt{{[}End dialogue{]}} \\
\\ 
\texttt{At the end of this dialogue, Harold would feel...} \\
\texttt{Disbelief} \\
\texttt{Irritation} \\
\texttt{Admiration} \\
\texttt{Indifference} \\
\\ 
\texttt{Give each of these possible emotions a score from 0{-}10 for the relative intensity that they are likely to be feeling each.} \\
\\ 
\texttt{You must output in the following format, including headings (of course, you should give your own scores), with no additional commentary:} \\
\\ 
\texttt{Disbelief: <score>} \\
\texttt{Irritation: <score>} \\
\texttt{Admiration: <score>} \\
\texttt{Indifference: <score>} \\
\\ 
\\ 
\texttt{{[}End of answer{]}} \\
\\ 
\texttt{Remember: zero is a valid score, meaning they are likely not feeling that emotion. You must score at least one emotion > 0.} \\
\\ 
\texttt{Your answer:} \\
            }
        }
    }
     \\ \\

    \theutterance \stepcounter{utterance}  
    & \multicolumn{4}{p{0.6\linewidth}}{
        \cellcolor[rgb]{0.561,0.737,0.561}{
            \makecell[{{p{\linewidth}}}]{
                \texttt{\tiny{[P1$\rangle$GM]}}
                \texttt{Disbelief: 7} \\
\texttt{Irritation: 6} \\
\texttt{Admiration: 0} \\
\texttt{Indifference: 0} \\
            }
        }
    }
    & & \\ \\

    \theutterance \stepcounter{utterance}  
    & & & \multicolumn{2}{p{0.3\linewidth}}{
        \cellcolor[rgb]{0.9,0.9,0.9}{
            \makecell[{{p{\linewidth}}}]{
                \texttt{\tiny{[GM$|$GM]}}
                \texttt{Disbelief: 7} \\
\texttt{Irritation: 6} \\
\texttt{Admiration: 0} \\
\texttt{Indifference: 0} \\
            }
        }
    }
    & & \\ \\

    \theutterance \stepcounter{utterance}  
    & & & \multicolumn{2}{p{0.3\linewidth}}{
        \cellcolor[rgb]{0.9,0.9,0.9}{
            \makecell[{{p{\linewidth}}}]{
                \texttt{\tiny{[GM$|$GM]}}
                \texttt{Disbelief: 4} \\
\texttt{Irritation: 6} \\
\texttt{Admiration: 0} \\
\texttt{Indifference: 0} \\
            }
        }
    }
    & & \\ \\

    \theutterance \stepcounter{utterance}  
    & & & \multicolumn{2}{p{0.3\linewidth}}{
        \cellcolor[rgb]{0.9,0.9,0.9}{
            \makecell[{{p{\linewidth}}}]{
                \texttt{\tiny{[GM$|$GM]}}
                \texttt{game\_result = LOSE} \\
            }
        }
    }
    & & \\ \\

    \theutterance \stepcounter{utterance}  
    & & \multicolumn{4}{p{0.6\linewidth}}{
        \cellcolor[rgb]{0.95,0.95,0.95}{
            \makecell[{{p{\linewidth}}}]{
                \texttt{\tiny{[P1$\leftrightarrow$P]}}
                \texttt{Post{-}episode Reflexion (not sent to GM):} \\
\texttt{New plan: I failed because I consistently over{-}indexed on "Irritation" (scoring it 8) while underestimating the cognitive dissonance that drives "Disbelief." Harold's opening statement ("It's not like the world is going to end") is a direct challenge to Lorelei's premise of existential stakes, which logically precedes emotional annoyance. His follow{-}up question ("And you think your candidate has all the answers?") reinforces this skepticism by questioning the efficacy of her candidate, not just her passion. While the condescending tone does generate irritation, the primary driver of his reaction is the logical gap between Lorelei's apocalyptic framing and his pragmatic view. Therefore, "Disbelief" should be the dominant emotion (score 9) as it captures the core of his argumentative stance, while "Irritation" remains a secondary, moderate reaction (score 4) to her intensity. "Admiration" and "Indifference" remain 0.} \\
\\ 
\texttt{Disbelief: 9} \\
\texttt{Irritation: 4} \\
\texttt{Admiration: 0} \\
\texttt{Indifference: 0} \\
            }
        }
    }
    & \\ \\

\end{supertabular}
}

\end{document}
